\section*{Discussion}

This report set out to test whether TabPFN's small-data advantage carries over to cancer survival prediction from multiomics. First, the direction of the original claim reproduced on TabPFN's own benchmark datasets under a simpler protocol: TabPFN had the highest mean ROC AUC among default-setting models, and its generations and TabFM were nearly interchangeable. Second, within single cancers, TabPFN remained competitive and was frequently the best model, but survival was hard for every method and no nonlinear model was significantly better than another. Third, pooling cancers raised every capable model by a similar margin (Fig.~\ref{fig:pooling}a), and controls without patient biology reached comparable scores (Fig.~\ref{fig:pooling}b). Fourth, pooled training lowered within-cancer ranking, more for the foundation models than for tree ensembles (Fig.~\ref{fig:effect}).

Our verdict is that TabPFN is reusable for the single-cohort, fixed-horizon survival tasks we studied. It runs without tuning, performed comparably to default-setting tree ensembles (a comparison that cannot establish equivalence), and had slightly lower log loss and Brier score than the tree ensembles in the pooled tasks (Extended Data Fig.~\ref{fig:ed_prob}), although neither measure establishes calibration.\citep{van2019calibration} This is in line with the independent clinical comparison, in which TabPFN's advantage over tuned established methods was inconsistent and generally within 0.02 AUROC.\citep{shaktah2026established} Our baselines were not tuned, and in large benchmarks tuning and ensembling can change rankings,\citep{erickson2025tabarena} so the comparison is between default settings. Its pooled pan-cancer scores, however, should not be read as patient-level prognosis.

We do not have a mechanistic account of why pooling hurt the foundation models more, and we offer one as a hypothesis to be tested. TabPFN attends across training samples and across features and fits no task-specific weights,\citep{hollmann2025accurate} and we hypothesise that its prediction for a patient therefore depends on the training rows most similar to that patient. We did not measure attention or run an ablation, so the architecture paper describes how the model computes and not why pooling hurt it. In the pooled table, the most similar rows are plausibly those from the same cohort, which the zero pattern identifies almost perfectly, and they share the cohort's average outcome. Tree ensembles build predictions from splits on individual features,\citep{breiman2001random,chen2016xgboost} which may make them less able to collapse onto a single latent variable. The hypothesis is directly testable. Standardizing features within each cancer and removing the zero fingerprint before pooling should restore within-cancer ranking for the foundation models if it is right, and should leave the gap unchanged if it is wrong; designs that add uninformative feature blocks or remove modalities have been used in a similar way to probe multiomics survival models.\citep{wissel2023systematic} This is the analogue of the remedies tested in previous reusability reports, and we have not yet run it.

Three bodies of work place the result in context without explaining it. In transfer learning, negative transfer is defined by comparison with learning on the target alone, which is the comparison in Fig.~\ref{fig:effect}, although that work concerns image domain adaptation and not in-context learning on pooled tables.\citep{wang2019characterizing} In prediction-model validation, discrimination depends on case mix, and an average over settings can conceal wide variation between them.\citep{dejong2023propensity,riley2016external} Our within-cancer scoring stratifies by cancer; it does not reweight patients to a target population, as propensity-based standardization does.\citep{dejong2023propensity} In routine health data, the presence or absence of a measurement can itself carry information, and models that exploit it can lose accuracy when the reason for missingness changes after deployment.\citep{sisk2021informative,rockenschaub2024robust} The zero pattern in our data records which assays were run for a cohort and not clinical behaviour, so these studies are analogies and not evidence about our mechanism. Finally, because TCGA molecular classes are organized largely by cell of origin,\citep{hoadley2018cell} the cohort information that models can use is partly biological, and the shortcut should not be dismissed as batch effect alone.

A reasonable counterargument is that, for an unselected pan-cancer population, cohort membership is part of the prediction problem. In oncology, though, the cancer type is known before any molecular profile is taken, so a pooled score answers a different question from the one asked at diagnosis. Within-cancer results answer the clinical question; pooled results, reported alone, can overstate it. Modest within-cancer discrimination is not unique to these models: in an 18-cohort TCGA benchmark, the best censored-survival method exceeded a clinical-variables-only reference by 0.002 to 0.028 in C-index in breast, head and neck, and both lung cancers, with the winning method chosen retrospectively in each cancer; C-index and ROC AUC are different measures and the values are not comparable.\citep{herrmann2021benchmark} The extreme-survival endpoint deserves the same caution. It produced the highest scores in the study but does not estimate performance for the fixed-horizon population, and its controls were the strongest of the three endpoints.

Using these models was straightforward. TabPFN and TabFM installed cleanly, needed no tuning, and ran on CPU. Three additions would make them easier to reuse in biomedical work. First, clear guidance on the 1,000-sample CPU safeguard of TabPFN v2 to v2.6, which excluded those versions from our largest task and is a runtime safeguard distinct from the models' recommended data sizes (for example 10,000 rows for TabPFN v2 and 50,000 for v2.5).\citep{hollmann2025accurate,priorlabs2025tabpfn25} Second, documented behavioral differences between generations, which our results suggest are small on public benchmarks but may matter under cohort structure. Third, a supported interface for censored time-to-event outcomes.

The limitations follow from the design. Fixed-horizon labels discard patients censored before the horizon, about half the cohort at 3 years. The reproduction used six datasets and one split each, raw ROC AUC and untuned baselines, whereas the original used 29 datasets, ten seeds and normalised scores, and we did not tabulate its per-dataset values. Clinical variables were excluded from our models, and survival benchmarks that include them use censored time-to-event metrics, so neither their results nor their conclusions about clinical-only references transfer directly to ours.\citep{herrmann2021benchmark,li2024combining} Cancer-specific test sets were small. The shortcut controls came from a single earlier split. The cancer-held-out diagnostic used a linear model with label-free feature selection that included the held-out cancer, and it was not repeated with foundation models; internal-external validation that leaves out one study at a time is the standard way to check transport across settings,\citep{riley2016external} and our diagnostic is a weaker version of it. The paired within-cancer analysis bootstraps patients, not model retraining. Finally, all data came from TCGA and CPTAC. The next steps are censored survival modeling applied identically to every model, the within-cancer standardization test described above, a comparison with clinical stage and age, and validation in an independent cohort. Until then, a pooled score is the beginning of an argument about biology, not its conclusion.
